\section{Kernel frameworks for device drivers}

\begin{frame}
  \frametitle{Kernel and Device Drivers}
  \begin{columns}
    \column{0.7\textwidth} In Linux, a driver is always interfacing
    with:
    \begin{itemize}
    \item a {\bf framework} that allows the driver to expose the
      hardware features to user space applications.
    \item a {\bf bus infrastructure}, part of the device model, to
      detect/communicate with the hardware.
    \end{itemize}
    This section focuses on the {\em kernel frameworks}, while the
    {\em bus infrastructure} was covered earlier in this training.
    \column{0.3\textwidth}
    \includegraphics[height=0.8\textheight]{slides/kernel-frameworks/driver-architecture.pdf}
  \end{columns}
\end{frame}

\subsection{User space vision of devices}

\begin{frame}
  \frametitle{Types of devices} Under Linux, there are essentially
  four types of devices:
  \begin{itemize}
  \item {\bf Network devices}. They are represented as network
    interfaces, visible in user space using \code{ip a}
  \item {\bf Block devices}. They are used to provide user space
    applications access to raw storage devices (hard disks, USB
    keys). They are visible to the applications as {\em device files}
    in \code{/dev}.
  \item {\bf Character devices}. They are used to provide user space
    applications access to all other types of devices (input, sound,
    graphics, serial, etc.). They are also visible to the applications
    as {\em device files} in \code{/dev}.
  \item {\bf Sysfs devices}. They don't have any of the above user space
    interfaces, only a representation in sysfs. "Internal" device drivers
    fall under this (e.g. pinctrl), but also some user-space accessible
    devices. E.g. gpio (deprecated), IIO (Industrial I/O).
  \end{itemize}
  $\rightarrow$ Most devices are {\em character devices}, so we will
  study these in more details.
\end{frame}

\begin{frame}
  \frametitle{Devices: everything is a file}
  \begin{itemize}
  \item A very important UNIX design decision was to represent most
    {\em system objects} as files
  \item It allows applications to manipulate all {\em system objects} with
    the normal file API (\code{open}, \code{read}, \code{write},
    \code{close}, etc.)
  \item So, devices had to be represented as files to the applications
  \item This is done through a special artifact called a {\bf device
      file}
  \item It is a special type of file, associating a file name
    visible to user space applications to a triplet the kernel understands:
    \begin{itemize}
    \item {\em type}: \code{char} or \code{block}
    \item {\em major}: class of the device
    \item {\em minor}: unique identifier in a class
    \end{itemize}
  \item All {\em device files} are by convention stored in the
    \code{/dev} directory
  \end{itemize}
\end{frame}

\begin{frame}[fragile]
\frametitle{Device files examples}

Example of device files in a Linux system

\small
\begin{verbatim}
$ ls -l /dev/ttyS0 /dev/tty1 /dev/sda /dev/sda1 /dev/sda2 /dev/sdc1 /dev/zero
brw-rw---- 1 root disk    8,  0 2011-05-27 08:56 /dev/sda
brw-rw---- 1 root disk    8,  1 2011-05-27 08:56 /dev/sda1
brw-rw---- 1 root disk    8,  2 2011-05-27 08:56 /dev/sda2
brw-rw---- 1 root disk    8, 32 2011-05-27 08:56 /dev/sdc
crw------- 1 root root    4,  1 2011-05-27 08:57 /dev/tty1
crw-rw---- 1 root dialout 4, 64 2011-05-27 08:56 /dev/ttyS0
crw-rw-rw- 1 root root    1,  5 2011-05-27 08:56 /dev/zero
\end{verbatim}
\normalsize

Example C code that uses the usual file API to write data to a serial port

\small
\begin{minted}{c}
int fd;
fd = open("/dev/ttyS0", O_RDWR);
write(fd, "Hello", 5);
close(fd);
\end{minted}
\end{frame}

\begin{frame}
  \frametitle{Creating device files}
  \begin{itemize}
    \item Before Linux 2.6.32, on basic Linux systems,
    the device files had to be created manually using the
    \code{mknod} command
    \begin{itemize}
    \item \code{mknod /dev/<device> [c|b] major minor}
    \item Needed root privileges
    \item Coherency between device files and devices handled by the
      kernel was left to the system developer
    \end{itemize}
  \item The \code{devtmpfs} virtual filesystem can be mounted on
    \code{/dev} and contains all the devices registered to kernel frameworks.
    The \kconfig{CONFIG_DEVTMPFS_MOUNT} kernel configuration option
    makes the kernel mount it automatically at boot time, except
    when booting on an initramfs.
  \item \code{devtmpfs} can be supplemented by userspace tools like
    \code{udev} or \code{mdev} to adjust permission/ownership, load
    kernel modules automatically and create symbolic links to devices.
  \end{itemize}
\end{frame}

\subsection{Character drivers}

\begin{frame}
  \frametitle{A character driver in the kernel}
  \begin{itemize}
  \item From the point of view of an application, a {\em character
      device} is essentially a {\bf file}.
  \item Character device drivers therefore implement {\bf operations}
    that let applications think the device is a file.
  \item In order to achieve this, a character driver implements
    the operations it wants from the \kstruct{file_operations}
    structure: \code{read}, \code{write}, \code{ioctl}, etc.
  \item The Linux filesystem layer will ensure that the driver's
    operations are called when a user space application makes the
    corresponding system call.
  \end{itemize}
\end{frame}

\begin{frame}
  \frametitle{From user space to the kernel: character devices}
  \begin{center}
    \includegraphics[height=0.8\textheight]{slides/kernel-frameworks/user-kernel-exchanges.pdf}
  \end{center}
\end{frame}

\begin{frame}[fragile]
  \frametitle{File operations}
  Here are the most important operations for a character
  driver, from the definition of \kstruct{file_operations}:
\begin{minted}[fontsize=\footnotesize]{c}
struct file_operations {
    struct module *owner;
    ssize_t (*read) (struct file *, char __user *,
        size_t, loff_t *);
    ssize_t (*write) (struct file *, const char __user *,
        size_t, loff_t *);
    long (*unlocked_ioctl) (struct file *, unsigned int,
        unsigned long);
    int (*mmap) (struct file *, struct vm_area_struct *);
    int (*open) (struct inode *, struct file *);
    int (*release) (struct inode *, struct file *);
    ...
};
\end{minted}
Many operations exist, they are all optional.
\end{frame}

\begin{frame}[fragile]
  \frametitle{open() and release()}
  \begin{itemize}
  \item \mint{c}+int foo_open(struct inode *i, struct file *f)+
    \begin{itemize}
    \item Called when user space opens the device file.
    \item {\bf Only implement this function when you do something
          special with the device at \code{open()} time.}
    \item \kstruct{inode} is a structure that uniquely represents a file
      in the filesystem (be it a regular file, a directory, a symbolic
      link, a character or block device)
    \item \kstruct{file} is a structure created every time a file is
      opened. Several file structures can point to the same
      \code{inode} structure.
      \begin{itemize}
      \item Contains information like the current position, the
        opening mode, etc.
      \item Has a \code{void *private_data} pointer that one can
        freely use.
      \item A pointer to the \ksym{file} structure is passed to all other
        operations
      \end{itemize}
    \end{itemize}
  \item \mint{c}+int foo_release(struct inode *i, struct file *f)+
    \begin{itemize}
    \item Called when user space closes the file.
    \item {\bf Only implement this function when you do something
          special with the device at \code{close()} time.}
    \end{itemize}
  \end{itemize}
\end{frame}

\begin{frame}[fragile]
  \frametitle{read() and write()}
  \begin{itemize}
  \item \mint[fontsize=\small]{c}+ssize_t foo_read(struct file *f, char __user *buf, size_t sz, loff_t *off)+
    \begin{itemize}
    \item Called when user space uses the \code{read()} system call on
      the device.
    \item Must read data from the device, write at most \code{sz}
      bytes to the user space buffer \code{buf}, and update the
      current position in the file \code{off}. \code{f} is a pointer
      to the same file structure that was passed in the \code{open()}
      operation
    \item Must return the number of bytes read.\\
	  \code{0} is usually interpreted by userspace as the end of
          the file.
    \item On UNIX, \code{read()} operations typically block when there
      isn't enough data to read from the device
    \end{itemize}
  \item \mint[fontsize=\small]{c}+ssize_t foo_write(struct file *f, const char __user *buf, size_t sz, loff_t *off)+
    \begin{itemize}
    \item Called when user space uses the \code{write()} system call
      on the device
    \item The opposite of \code{read}, must read at most \code{sz}
      bytes from \code{buf}, write it to the device, update \code{off}
      and return the number of bytes written.
    \end{itemize}
  \end{itemize}
\end{frame}

\begin{frame}
  \frametitle{Exchanging data with user space 1/3}
  \begin{itemize}
  \item Kernel code isn't allowed to directly access user space
    memory, using \kfunc{memcpy} or direct pointer dereferencing
    \begin{itemize}
    \item User pointer dereferencing is disabled by default to make it
      harder to exploit vulnerabilities.
    \item If the address passed by the application was invalid, the
      kernel could segfault.
    \item {\bf Never} trust user space. A malicious application could
      pass a kernel address which you could overwrite with device data
      (\code{read} case), or which you could dump to the device
      (\code{write} case).
    \item Doing so does not work on some architectures anyway.
    \end{itemize}
  \item To keep the kernel code portable, secure, and have proper
    error handling, your driver must use special kernel functions
    to exchange data with user space.
  \end{itemize}
\end{frame}

\begin{frame}[fragile]
  \frametitle{Exchanging data with user space 2/3}
  \begin{itemize}
  \item A single value
    \begin{itemize}
    \item \code{get_user(v, p);}
      \begin{itemize}
      \item The kernel variable \code{v} gets the value pointed by the
        user space pointer \code{p}
      \end{itemize}
    \item \code{put_user(v, p);}
      \begin{itemize}
      \item The value pointed by the user space pointer \code{p} is
        set to the contents of the kernel variable \code{v}.
      \end{itemize}
    \end{itemize}
  \item A buffer
    \begin{itemize}
    \item
      \begin{minted}{c}
unsigned long copy_to_user(void __user *to, const void *from,
                           unsigned long n);
      \end{minted}
    \item
      \begin{minted}{c}
unsigned long copy_from_user(void *to, const void __user *from,
                             unsigned long n);
      \end{minted}
    \end{itemize}
  \item The return value must be checked. Zero on success, non-zero on
    failure. If non-zero, the convention is to return \code{-}\ksym{EFAULT}.
  \end{itemize}
\end{frame}

\begin{frame}
 \frametitle{Exchanging data with user space 3/3}
 \begin{center}
    \includegraphics[height=0.8\textheight]{slides/kernel-frameworks/copy-to-from-user.pdf}
 \end{center}
\end{frame}

\begin{frame}
  \frametitle{Zero copy access to user memory}
  \begin{itemize}
  \item Having to copy data to or from an intermediate kernel buffer
    can become expensive when the amount of data to transfer is
    large (video).
  \item \emph{Zero copy} options are possible:
    \begin{itemize}
    \item \code{mmap()} system call to allow user space to directly
      access memory mapped I/O space. See our \code{mmap()} chapter.
    \item \kfunc{get_user_pages} and related functions to get a
      mapping to user pages without having to copy them.
    \end{itemize}
  \end{itemize}
\end{frame}

\begin{frame}[fragile]
  \frametitle{unlocked\_ioctl()}
  \begin{itemize}
  \item \mint[fontsize=\small]{c}+long unlocked_ioctl(struct file *f, unsigned int cmd, unsigned long arg)+
    \begin{itemize}
    \item Associated to the \code{ioctl()} system call.
    \item Called unlocked because it didn't hold the Big Kernel Lock
      (gone now).
    \item Allows to extend the driver capabilities beyond the limited
      read/write API.
    \item For example: changing the speed of a serial port, setting
      video output format, querying a device serial number... Used
      extensively in the V4L2 (video) and ALSA (sound) driver frameworks.
    \item \code{cmd} is a number identifying the operation to perform.\\
      See \kdochtml{driver-api/ioctl} for the recommended way of choosing
      \code{cmd} numbers.
    \item \code{arg} is the optional argument passed as third argument
      of the \code{ioctl()} system call. Can be an integer, an
      address, etc.
    \item The semantic of \code{cmd} and \code{arg} is
      driver-specific.
    \end{itemize}
  \end{itemize}
\end{frame}

\begin{frame}[fragile]
  \frametitle{ioctl() example: kernel side}
\begin{minted}[fontsize=\tiny]{c}
#include <linux/phantom.h>

static long phantom_ioctl(struct file *file, unsigned int cmd,
    unsigned long arg)
{
    struct phm_reg r;
    void __user *argp = (void __user *)arg;

    switch (cmd) {
    case PHN_SET_REG:
        if (copy_from_user(&r, argp, sizeof(r)))
            return -EFAULT;
        /* Do something */
        break;
    ...
    case PHN_GET_REG:
        if (copy_to_user(argp, &r, sizeof(r)))
            return -EFAULT;
        /* Do something */
        break;
    ...
    default:
        return -ENOTTY;
    }

    return 0;
}
\end{minted}
\small Selected excerpt from \kfile{drivers/misc/phantom.c}
\end{frame}

\begin{frame}[fragile]
  \frametitle{Ioctl() Example: Application Side}
\begin{minted}[fontsize=\footnotesize]{c}
#include <linux/phantom.h>

int main(void)
{
    int fd, ret;
    struct phm_reg reg;

    fd = open("/dev/phantom");
    assert(fd > 0);

    reg.field1 = 42;
    reg.field2 = 67;

    ret = ioctl(fd, PHN_SET_REG, &reg);
    assert(ret == 0);

    return 0;
}
\end{minted}
\end{frame}

\subsection{The concept of kernel frameworks}

\begin{frame}
  \frametitle{Beyond character drivers: kernel frameworks}
  \begin{itemize}
  \item Many device drivers are not implemented directly as character
    drivers
  \item They are implemented under a \emph{framework}, specific to a
    given device type (framebuffer, V4L, serial, etc.)
    \begin{itemize}
    \item The driver plugs into a framework API rather than into the
      generic \kstruct{file_operations}
    \item The framework implements once \kstruct{file_operations} to
      expose character devices to userspace
    \item That implementation is shared across all framework devices: it
      minimizes driver boilerplate and it provides a coherent userspace
      interface whatever driver is being used
    \end{itemize}
  \end{itemize}
\end{frame}

\begin{frame}
  \frametitle{Example: Some Kernel Frameworks}
  \begin{center}
    \includegraphics[height=0.8\textheight]{slides/kernel-frameworks/frameworks.pdf}
   \end{center}
 \end{frame}

\subsection{Example: the input subsystem}
  \begin{frame}
    \frametitle{What is the input subsystem?}
  \begin{itemize}
  \item The input subsystem takes care of all the input events coming
    from the human user.
  \item Initially written to support the USB {\em HID} (Human
    Interface Device) devices, it quickly grew up to handle all kinds
    of inputs (using USB or not): keyboards, mice, joysticks,
    touchscreens, etc.
  \item The input subsystem is split in two parts:
    \begin{itemize}
    \item {\bf Device drivers}: they talk to the hardware (for example
      via USB), and provide events (keystrokes, mouse movements,
      touchscreen coordinates) to the input core
    \item {\bf Event handlers}: they get events from drivers and pass
      them where needed via various interfaces (most of the time
      through \code{evdev})
    \end{itemize}
  \item In user space it is usually used by the graphic stack such
    as {\em X.Org}, {\em Wayland} or {\em Android's InputManager}.
  \end{itemize}
\end{frame}

\begin{frame}{Input subsystem diagram}
  \begin{center}
    \includegraphics[height=0.8\textheight]{slides/kernel-frameworks/input-subsystem-diagram.pdf}
  \end{center}
\end{frame}

\begin{frame}{Input subsystem overview}
  \begin{itemize}
  \item Kernel option \kconfig{CONFIG_INPUT}
    \begin{itemize}
    \item \code{menuconfig INPUT}
      \begin{itemize}
      \item \code{tristate "Generic input layer (needed for keyboard, mouse, ...)"}
      \end{itemize}
    \end{itemize}
  \item Implemented in \kdir{drivers/input}
    \begin{itemize}
    \item \krelfile{drivers/input}{input.c},
          \krelfile{drivers/input}{input-poller.c}, \krelfile{drivers/input}{evdev.c}...
    \end{itemize}
  \item Defines the user/kernel API
    \begin{itemize}
    \item \kfile{include/uapi/linux/input.h}
    \end{itemize}
  \item Defines the set of operations an input driver must implement
    and helper functions for the drivers
    \begin{itemize}
    \item \kstruct{input_dev} for the device driver part
    \item \kstruct{input_handler} for the event handler part
    \item  \kfile{include/linux/input.h}
    \end{itemize}
  \end{itemize}
\end{frame}

\begin{frame}[fragile]{Input subsystem API 1/3}
  An {\em input device} is described by a very long
  \kstruct{input_dev} structure, an excerpt is:
  \begin{block}{}
  \begin{minted}[fontsize=\tiny]{c}
struct input_dev {
    const char *name;
    [...]
    struct input_id id;
    [...]
    unsigned long evbit[BITS_TO_LONGS(EV_CNT)];
    unsigned long keybit[BITS_TO_LONGS(KEY_CNT)];
    [...]
    int (*getkeycode)(struct input_dev *dev,
                      struct input_keymap_entry *ke);
    [...]
    int (*open)(struct input_dev *dev);
    [...]
    int (*event)(struct input_dev *dev, unsigned int type,
                 unsigned int code, int value);
    [...]
};
\end{minted}
\end{block}
  Before being used, this structure must be allocated and
  initialized, typically with:
  \code{struct input_dev *devm_input_allocate_device(struct device *dev);}
\end{frame}

\begin{frame}[fragile]{Input subsystem API 2/3}
  \begin{itemize}
  \item Depending on the type of events that will be generated, the
    input bit fields \code{evbit} and \code{keybit} must be configured:
    For example, for a button we only generate
    \ksym{EV_KEY} type events, and from these only \ksym{BTN_0} events
    code:
    \begin{block}{}
    \begin{minted}[fontsize=\footnotesize]{c}
set_bit(EV_KEY, myinput_dev.evbit);
set_bit(BTN_0, myinput_dev.keybit);
    \end{minted}
    \end{block}
  \item Once the {\em input device} is allocated and filled, the
    function to register it
    is: \code{int input_register_device(struct input_dev *);}
  \end{itemize}
\end{frame}

\begin{frame}[fragile]{Input subsystem API 3/3}
  The events are sent by the driver to the event handler using
  \begin{minted}[fontsize=\footnotesize]{c}
void input_event(struct input_dev *dev, unsigned int type, unsigned int code, int value)
  \end{minted}
  \begin{itemize}
  \item The event types are documented in \kdochtml{input/event-codes}
  \item An event is composed by one or several input data changes
    (packet of input data changes) such as the button state, the
    relative or absolute position along an axis, etc..
  \item The input subsystem provides other wrappers such as:
    \begin{itemize}
    \item \kfunc{input_report_key}
    \item \kfunc{input_report_abs}
    \end{itemize}
  \end{itemize}
  After submitting potentially multiple events, the {\em input} core must be
  notified by calling:
  \begin{minted}[fontsize=\footnotesize]{c}
void input_sync(struct input_dev *dev)
  \end{minted}
\end{frame}

\begin{frame}[fragile]{Example from drivers/hid/usbhid/usbmouse.c}
  \begin{block}{}
  \begin{minted}[fontsize=\scriptsize]{c}
static void usb_mouse_irq(struct urb *urb)
{
        struct usb_mouse *mouse = urb->context;
        signed char *data = mouse->data;
        struct input_dev *dev = mouse->dev;
        ...

        input_report_key(dev, BTN_LEFT,   data[0] & 0x01);
        input_report_key(dev, BTN_RIGHT,  data[0] & 0x02);
        input_report_key(dev, BTN_MIDDLE, data[0] & 0x04);
        input_report_key(dev, BTN_SIDE,   data[0] & 0x08);
        input_report_key(dev, BTN_EXTRA,  data[0] & 0x10);

        input_report_rel(dev, REL_X,     data[1]);
        input_report_rel(dev, REL_Y,     data[2]);
        input_report_rel(dev, REL_WHEEL, data[3]);

        input_sync(dev);
        ...
}
  \end{minted}
  \end{block}
\end{frame}

\begin{frame}[fragile]{Polling input devices}
  \begin{itemize}
  \item The input subsystem provides an API to support simple input
    devices that {\em do not raise interrupts} but have to be {\em
      periodically scanned or polled} to detect changes in their
    state.
  \item Setting up polling is done using \kfunc{input_setup_polling}:
  \begin{minted}[fontsize=\footnotesize]{c}
int input_setup_polling(struct input_dev *dev, void (*poll_fn)(struct input_dev *dev));
  \end{minted}
  \item \code{poll_fn} is the function that will be called
    periodically.
  \item The polling interval can be set using
    \kfunc{input_set_poll_interval} or
    \kfunc{input_set_min_poll_interval} and
    \kfunc{input_set_max_poll_interval}
  \end{itemize}
\end{frame}

\begin{frame}[fragile]{{\em evdev} user space interface}
  \begin{itemize}
  \item The main user space interface to {\em input devices} is the
    {\bf event interface}
  \item Each {\em input device} is represented as a
    \code{/dev/input/event<X>} character device
  \item A user space application can use blocking and non-blocking
    reads, but also \code{select()} (to get notified of events) after
    opening this device.
  \item Each read will return \kstruct{input_event} structures of the
    following format:
    \begin{block}{}
    \begin{minted}[fontsize=\tiny]{c}
struct input_event {
        struct timeval time;
        unsigned short type;
        unsigned short code;
        unsigned int value;
};
\end{minted}
\end{block}
\item A very useful application for {\em input device} testing is
  \code{evtest}, from \url{https://cgit.freedesktop.org/evtest/}
  \end{itemize}
\end{frame}

\subsection{Device-managed allocations}

\begin{frame}
  \frametitle{Device managed allocations}
  \begin{itemize}
  \item The \code{probe()} function is typically responsible for
    allocating a significant number of resources: memory, mapping I/O
    registers, registering interrupt handlers, etc.
  \item These resource allocations have to be properly freed:
    \begin{itemize}
    \item In the \code{probe()} function, in case of failure
    \item In the \code{remove()} function
    \end{itemize}
  \item This required a lot of failure handling code that was rarely
    tested
  \item To solve this problem, {\em device managed} allocations have
    been introduced.
  \item The idea is to associate resource allocation with the
    \code{struct device}, and automatically release those resources
    \begin{itemize}
    \item When the device disappears
    \item When the device is unbound from the driver
    \end{itemize}
  \item Functions prefixed by \code{devm_}
  \item See \kdochtml{driver-api/driver-model/devres} for details
  \end{itemize}
\end{frame}

\begin{frame}[fragile]
  \frametitle{Device managed allocations: memory allocation example}
  \begin{itemize}
  \item Normally done with \code{kmalloc(size_t, gfp_t)}, released
    with \code{kfree(void *)}
  \item Device managed with \code{devm_kmalloc(struct device *,
      size_t, gfp_t)}
  \end{itemize}

  \begin{columns}
    \column[t]{0.5\textwidth}
    \begin{block}{Without devm functions}
\fontsize{6}{6}\selectfont
    \begin{minted}{c}
int foo_probe(struct platform_device *pdev)
{
        struct foo_t *foo = kmalloc(sizeof(struct foo_t),
                                    GFP_KERNEL);
        /* Register to framework, store
         * reference to framework structure in foo */
        ...
        if (failure) {
                kfree(foo);
                return -EBUSY;
        }
        platform_set_drvdata(pdev, foo);
        return 0;
}

void foo_remove(struct platform_device *pdev)
{
        struct foo_t *foo = platform_get_drvdata(pdev);
        /* Retrieve framework structure from foo
           and unregister it */
        ...
        kfree(foo);
}
    \end{minted}
  \end{block}
    \column[t]{0.5\textwidth}
    \begin{block}{With devm functions}
\fontsize{6}{6}\selectfont
    \begin{minted}{c}
int foo_probe(struct platform_device *pdev)
{
        struct foo_t *foo = devm_kmalloc(&pdev->dev,
                           sizeof(struct foo_t),
                           GFP_KERNEL);
        /* Register to framework, store
         * reference to framework structure in foo */
        ...
        if (failure)
                return -EBUSY;
        platform_set_drvdata(pdev, foo);
        return 0;
}

void foo_remove(struct platform_device *pdev)
{
        struct foo_t *foo = platform_get_drvdata(pdev);
        /* Retrieve framework structure from foo
           and unregister it */
        ...
        /* foo automatically freed */
}
    \end{minted}
  \end{block}
\end{columns}

\end{frame}

\begin{frame}
  \frametitle{Device managed allocations caveats}
  \begin{itemize}
  \item Cleanup is done when the \code{struct device} is cleaned up.
    There is no reference counting or anything like that.
  \item Don't use if the allocated memory is used outside of the device
    node. E.g. if the userspace device file is still open after remove.
  \item Be very careful when there are circular references.
  \item \href{https://lpc.events/event/16/contributions/1227/}{"Why is
    \code{devm_kzalloc()} harmful and what can we do about it",
    Laurent Pinchart, LPC 2022}
  \end{itemize}
\end{frame}


\subsection{Driver data structures and links}

\begin{frame}
  \frametitle{Driver data layout}
  Three main data structures:
  \begin{itemize}

  \item Bus-specific device structure (\kstruct{i2c_client},
    \kstruct{usb_dev}, etc)
  \begin{itemize}
  \item It \emph{always embeds} a \kstruct{device}
  \item A pointer to it is passed to the \code{probe()} function
  \end{itemize}

  \item Framework-specific device struct (\kstruct{input_dev},
    \kstruct{rtc_device}, etc)
  \begin{itemize}
  \item It \emph{might embed} a \kstruct{device} if the framework wants
    a sysfs device
  \item Careful! We might have one bus \kstruct{device} and one
    framework \kstruct{device} for the same device!
  \end{itemize}

  \item Driver private data
  \begin{itemize}
  \item The structure is driver specific, with space for all device
    state information
  \item It stores references to both the bus and framework devices
  \end{itemize}

  \end{itemize}
\end{frame}

\begin{frame}[fragile]
  \frametitle{Driver data allocation stategies}
  \begin{columns}
    \column{0.33\textwidth}
      \small Private data embeds the framework device therefore a single
      allocation allocates both the framework device and the private
      data.

      \begin{minted}[fontsize=\scriptsize]{c}
struct imx_port {
    struct uart_port port;
    struct timer_list timer;
    unsigned int old_status;
    int txirq, rxirq, rtsirq;
    [...]
};

sport = devm_kzalloc(&pdev->dev,
                sizeof(*sport),
                GFP_KERNEL);
if (!sport)
    return -ENOMEM;
      \end{minted}

    \column{0.33\textwidth}
      \small The framework exposes an helper to allocate the framework device,
      with space at the end to put the private data.

      \begin{minted}[fontsize=\scriptsize]{c}
struct da311_data *data;
struct iio_dev *idev;

idev = devm_iio_device_alloc(
                &client->dev,
                sizeof(*data));
if (!idev)
  return -ENOMEM;

data = iio_priv(idev);
data->client = client;
      \end{minted}

    \column{0.33\textwidth}
      \small The framework device and private data are allocated separately.

      \begin{minted}[fontsize=\scriptsize]{c}
struct rtc_device *rtc;
struct ds1305 *ds1305;

ds1305 = devm_kzalloc(&spi->dev,
                sizeof(*ds1305),
                GFP_KERNEL);
if (!ds1305)
    return -ENOMEM;

rtc = devm_rtc_allocate_device(
                &spi->dev);
if (IS_ERR(rtc))
    return PTR_ERR(rtc);
      \end{minted}
  \end{columns}
\end{frame}

\begin{frame}
  \frametitle{Links between data structures}

  \begin{columns}

  \column{0.6\textwidth}

    \begin{itemize}
    \item Inside bus callbacks, we get passed our bus device
      \begin{itemize}
      \item Inside \kstruct{device}, a pointer-sized field
        \code{dev->driver_data} is reserved for driver usage
      \item Use \kfunc{dev_set_drvdata} at \code{probe()} time to put a
        reference to our private data
      \item From bus callbacks, we can retrieve our private
        data using \kfunc{dev_get_drvdata}
      \end{itemize}

    \item Inside framework callbacks, we get passed our framework device
      \begin{itemize}
      \item If our framework device is embedded in our private data, we
        use \kfunc{container_of} that works using compiler provided
        \code{offsetof()}
      \item Otherwise, we use the framework device \code{dev->driver_data}
        and retrieve our private data reference
      \end{itemize}
    \end{itemize}

  \column{0.4\textwidth}
    \includegraphics[height=0.8\textheight]{slides/kernel-frameworks/link-structures.pdf}
  \end{columns}
\end{frame}

\setuplabframe
{Expose the Nunchuk to user space}
{
  \begin{itemize}
  \item Extend the Nunchuk driver to expose the Nunchuk features to
    user space applications, as an {\em input} device.
  \item Test the operation of the Nunchuk using \code{evtest}
  \end{itemize}
}
